\section{Results}
\label{sec:experiments}
\begin{table}[t]
\caption{GeoToolEval baseline comparison. AgentGS process metrics are per-record averages from the present 603-record end-to-end experiment; FAS values are retained from the source manuscript. Direct-answer language models have no tool traces; their process metrics are not applicable (---).}
\label{tab:main-results}
\centering
\begin{tabularx}{\columnwidth}{@{}Xrrr@{}}
\toprule
Method & FAS & TSF1 & TOS \\
\midrule
AgentGS (GPT-5.6 Terra) & \textbf{0.8731} & 0.9170 & 0.9166 \\
GeoAgent (GPT-5.6 Terra) & 0.7804 & 0.8663 & 0.7911 \\
EarthAgent (GPT-5.6 Terra) & 0.0862 & 0.1182 & 0.1390 \\
DeepSeek V4 Pro & 0.5162 & --- & --- \\
GPT-5.6 Terra & 0.3137 & --- & --- \\
GPT-5.4 & 0.3069 & --- & --- \\
Qwen-Plus & 0.2918 & --- & --- \\
GeoGPT & 0.2500 & --- & --- \\
\bottomrule
\end{tabularx}
\end{table}

We evaluate AgentGS on GeoToolEval and GeoBenchX, examining answer quality, tool use, and performance across workflow lengths.

\emph{Draft note:} Numerical comparisons below retain the previous evaluation (Appendix~\ref{app:experimental-details}); the expanded comparison in Table~\ref{tab:main-results}, including Quality, PAS, and ECR, remains a placeholder pending reruns.

\subsection{Experimental Setup}
\label{sec:experimental-setup}
The GeoToolEval evaluation covers all 603 records, with failures retained in the averages. We report final-answer agreement and tool-workflow metrics defined in Section~\ref{sec:metrics}. The expanded comparison includes eight AgentGS backbones and Bash-only agents using DeepSeek V4.1 Flash and GPT-5.6 Terra (Table~\ref{tab:main-results}). Bash-only agents explore the data through shell commands and are compared by FAS.

The previous evaluation uses GPT-5.6 Terra for AgentGS, GeoAgent~\cite{chen2024geoagent}, and EarthAgent~\cite{feng2026earthagent}, together with direct-answer language-model baselines. Those direct-answer results are distinct from the planned Bash-only comparison. Their scores are retained in Table~\ref{tab:prior-results}; run provenance and configuration details appear in Appendix~\ref{app:experimental-details}.

\subsection{Main Results on GeoToolEval}
\label{sec:main-results}
AgentGS achieves the highest mean FAS among the previously evaluated methods, scoring 0.8731 compared with 0.7804 for GeoAgent and 0.0862 for EarthAgent (Table~\ref{tab:prior-results}). Its advantage over the strongest agent baseline is 9.27 percentage points. Among direct-answer models, DeepSeek V4 Pro reaches 0.5162, while the same GPT-5.6 Terra backbone without tool execution scores 0.3137. These results support the effectiveness of AgentGS for completing the benchmark's scientific analyses.

The answer advantage is accompanied by higher agreement with reference tool workflows. AgentGS obtains TSF1/TOS of 0.9170/\allowbreak 0.9166, compared with 0.8663/\allowbreak 0.7911 for GeoAgent and 0.1182/\allowbreak 0.1390 for EarthAgent. Relative to GeoAgent, the gains are 5.07 percentage points in tool selection and 12.55 points in relative order. The larger order-score gain is consistent with the role of structured dependencies in organizing execution; Section~\ref{sec:ablations} examines individual mechanisms.

Figure~\ref{fig:geotooleval}b presents domain-specific FAS. This breakdown complements the overall score because themes differ in size, data type, and workflow composition. The reported aggregate gains do not imply an equal benefit across all domains.

\subsection{Performance across Workflow Lengths}
AgentGS obtains FAS of 0.9425 on single-step tasks, 0.7557 on tasks with two or three reference calls, and 0.6167 on tasks with at least four calls (Table~\ref{tab:chain-results}). Performance decreases as the evaluated workflows become longer.
\begin{table}[t]
\caption{AgentGS by reference-chain length on GeoToolEval. Process scores for the single-step stratum are not supplied in the reported breakdown (---).}
\label{tab:chain-results}
\centering
\begin{tabular}{@{}lrrrr@{}}
\toprule
Reference calls & Records & FAS & TSF1 & TOS \\
\midrule
1 & 424 & 0.9425 & --- & --- \\
2--3 & 119 & 0.7557 & 0.9779 & 0.9832 \\
$\geq4$ & 60 & 0.6167 & 0.8013 & 0.7785 \\
\bottomrule
\end{tabular}
\end{table}


On the 119 two- to three-step records, TSF1/TOS reach 0.9779/\allowbreak 0.9832. On the 60 longer records, they decrease to 0.8013/\allowbreak 0.7785. Longer workflows therefore remain challenging in both answer quality and tool orchestration. High process agreement on shorter workflows also coexists with lower FAS, showing that selecting and ordering the expected tools alone does not ensure a correct answer. Domain and length effects overlap: remote sensing accounts for 42 of the 60 longest-chain records.

\subsection{External Evaluation on GeoBenchX}
\label{sec:external-evaluation}
We further evaluate AgentGS on GeoBenchX~\cite{krechetova2025geobenchx} by integrating its official tools and constructing the corresponding Skills and operation mappings. AgentGS and the official ReAct baseline are each reported on 202 tasks using GPT-5.6 Terra. The same model judges reference-solution agreement on a 0--2 scale; we report mean score and the rate of full matches.
\begin{table}[t]
\caption{Evaluation results on GeoBenchX. Both methods use GPT-5.6 Terra.}
\label{tab:external-results}
\centering
\begin{tabularx}{\columnwidth}{@{}Xrr@{}}
\toprule
Method & Mean score & Full-match \\
 & $(0$--$2)\uparrow$ & rate $(\%)\uparrow$ \\
\midrule
ReAct & 1.1980 & 53.47 \\
AgentGS & \textbf{1.3762} & \textbf{62.87} \\
\bottomrule
\end{tabularx}
\end{table}


AgentGS achieves a mean score of 1.3762, compared with 1.1980 for ReAct, and a full-match rate of 62.87\%, compared with 53.47\%. The improvements are 0.1782 points and 9.40 percentage points, respectively. Both measures favor AgentGS in this additional tool environment. This evaluation uses dedicated strategy knowledge, and verification of paired task membership and matched conditions remains necessary for a controlled transfer claim (Appendix~\ref{app:experimental-details}).
